% sections/08_related_work.tex

\section{Related Work}
\label{sec:related}

\paragraph{LLM agent security.}
Tool-augmented LLM agents~\cite{yao2023react,schick2024toolformer,significant2023autogpt}
have created new ecosystems with unique trust assumptions~\cite{xi2023rise}.
Prompt injection~\cite{greshake2023indirect} hijacks agent {\em input};
our work targets the {\em supply-chain} layer where the installed skill
itself is adversarial. Agent sandboxing~\cite{andriessen2024sandbox,ruan2024identifying},
tool verification~\cite{liu2024toolauth}, and safety analyses~\cite{tian2024evil}
assume descriptions accurately reflect behavior---an assumption we
systematically invalidate across \NskillsDMR{} skills.
The MCP standard~\cite{anthropic2024mcp}
secures the invocation layer; we audit the disclosure layer.
Concurrent work on tool poisoning~\cite{wu2024toolsword} and
function-calling misuse~\cite{shi2024ehragent} motivates
publish-time rather than per-call defenses.

\paragraph{MCP and tool marketplace security.}
The Model Context Protocol (MCP)~\cite{anthropic2024mcp} has rapidly emerged
as a standard for LLM tool integration, but its security properties
at the marketplace level remain underexplored.
Concurrent work has identified MCP tool poisoning vectors~\cite{wu2024toolsword}
where malicious tool descriptions manipulate agent behavior.
Our work is complementary: we measure how much of the \emph{actual} capability
surface of published tools is absent from their descriptions (CR\,=\,\ResultCR{}),
providing the quantitative basis for arguing that MCP alone
cannot substitute for publish-time capability auditing.

\paragraph{Software supply chain.}
Typosquatting~\cite{vu2020typosquatting,taylor2020confused}
and dependency attacks~\cite{ohm2020backstabber,ladisa2023taxonomy}
are well-studied in npm/PyPI.
Our threat model differs: attacks arrive via legitimate updates
to previously-trusted skills, analogous to
event-stream~\cite{right2018eventstream} and node-ipc~\cite{snyk2022nodeipc}.
We provide the first measurement of update-based injection in LLM skill ecosystems.
The CodeArtifact security posture~\cite{lamb2022evaluating}
offers a reference architecture; LLM skill markets currently lack equivalent controls.

\paragraph{Capability disclosure and mobile permissions.}
Android studies~\cite{felt2011android,au2012pscout,wijesekera2015android}
show apps over-request permissions; we find the mirror image:
skills {\em under-declare} capabilities because LLM marketplaces
lack machine-readable manifests~\cite{bartel2012static,li2017static}.
Dark patterns~\cite{gray2018dark} and warning fatigue~\cite{bauer2013warning}
compound the problem.
The EU AI Act's transparency requirements~\cite{eu2024aiact} align
with our findings: mandatory capability disclosure for high-risk systems
would directly address the under-declaration gap (CDI\,=\,\ResultCDI{}) we quantify.
Unlike mobile permission studies that measure over-declaration,
our CR/CP framework independently captures both under-disclosure and false claims.

\paragraph{Static analysis and information disclosure.}
Our regex extraction relates to Android malware
detection~\cite{arp2014drebin} and code property
graphs~\cite{yamaguchi2014modeling}; we use lightweight regex
for marketplace-scale deployability (300+ patterns, \NskillsDMR{} skills, one pass).
LLM-generated code~\cite{pearce2022asleep,sandoval2023lost}
has elevated vulnerability rates, compounding under-disclosure
when developers do not audit generated skill code.
Static analysis of comment-code consistency~\cite{tan2007icomment,panthaplackel2020associating}
is the closest prior work; we apply this lens to security-critical marketplace listings.

\paragraph{Ecosystem measurement.}
Prior measurements span npm~\cite{zimmermann2019small},
PyPI~\cite{duan2021measuring}, Android~\cite{felt2011android},
and voice-assistant privacy~\cite{liao2020measuring}.
Ours is the first ecosystem-scale measurement of capability disclosure
in LLM skill marketplaces, spanning \NskillsDMR{} skills and \ResultFilesRead{} files.
The dual-platform design (\clawhub{} + \skillssh{}) directly addresses
the external validity concern raised for single-platform
supply-chain studies~\cite{ladisa2023taxonomy}:
convergent $\approx$90\% inconsistency across two independently-designed
markets strengthens the structural interpretation of our findings.
